\section*{Agradecimentos}

Ao Prof.\ Alexandre, pela supervisão e pela liberdade de deixar o projeto crescer além do escopo inicial.

À Maria Licia de Lima Farias e ao Prof.\ Dr.\ Paulo Eduardo de Brito, cujo trabalho de conclusão de curso \emph{Copérnico estava certo? Uma análise de epiciclos e séries de Fourier} \cite{farias2016} deu a ideia que originou este projeto. A pergunta que aquele trabalho fez, se uma soma de círculos girando consegue descrever a órbita de um planeta, é a pergunta de que este relatório parte. Os dados históricos dos epiciclos de Copérnico que usamos no \autoref{sec:copernico-revisitado} foram organizados por ele.

Ao Jet Propulsion Laboratory (JPL/NASA), pelo serviço Horizons e pelas efemérides DE440/DE441; ao USGS Astrogeology Science Center e ao Solar System Scope, pelos mapas de superfície; e aos mantenedores do \emph{Yale Bright Star Catalogue}.
Parte da programação, da revisão de código e da organização deste texto contou com o apoio do assistente de inteligência artificial Claude, da Anthropic, usado como ferramenta de desenvolvimento. As decisões de projeto, a verificação dos resultados e a responsabilidade pelo conteúdo são dos autores.


